The following tables summarize the ring versions discussed in this book, rather than every variant of these scheme families. Set $R=\mathbb Z[X]/(X^n+1)$ and $R_q=R/qR$, with power-of-two $n$. Write $[x]_q$ for a centered representative and $\lfloor x\rceil$ for nearest-integer rounding; Rounding, modular reduction, and coefficient bounds are interpreted coefficientwise; polynomial multiplication is in the stated quotient ring.

We use the plus-phase convention $b+\langle\vec a,\vec s\rangle$ or $B+AS$ throughout these tables. The TFHE chapter uses the equivalent minus convention, obtained by negating the public mask; blind-rotation and key-switching signs must follow that choice. Evaluated ciphertext errors are not assumed to retain the distribution or standard deviation of fresh encryption. A summary formula is subject to its decoding bound, compatible parameters, and evaluation-key assumptions. In particular, publishing encryptions of a secret under a related secret may require a circular/KDM-security assumption or an appropriate key hierarchy.

This requirement depends on the actual auxiliary-message function and ciphertext format. Same-key raw $R_q$-affine phases admit a mask-shift reduction; that result does not automatically cover secret squares or arbitrary coefficient maps. Independently generated keys in a directed acyclic graph admit a suitable hybrid argument even for such efficiently computable source-key functions, subject to the stated sampling and dependency conditions. Neither proof automatically establishes security for a fixed public-seed expansion. See Section~\ref{sec:glwe-key-switching} for the exact conditions.

\captionsetup{justification=raggedright,singlelinecheck=false,labelfont=bf}


\begin{table}[htbp]
\small
\renewcommand{\arraystretch}{1.15}
\begin{tabularx}{\linewidth}{|p{0.17\linewidth}|X|}
\hline
\textbf{TFHE}&LWE for scalar ciphertexts; GLWE/RLWE for ring ciphertexts and evaluation keys.\\\hline
\textbf{BFV / CKKS / BGV}&RLWE for the ring versions discussed here. Security depends on the ring, dimensions, moduli, secret/error distributions, and released evaluation keys.\\\hline
\end{tabularx}
\caption{\textbf{Hard Problem Basis}}
\end{table}

\begin{table}[htbp]
\small
\renewcommand{\arraystretch}{1.15}
\begin{tabularx}{\linewidth}{|p{0.17\linewidth}|X|}
\hline
\textbf{TFHE}&Scalar interface: $(\vec a,b)\in\mathbb Z_q^{k+1}$. Blind rotation uses GLWE/GGSW ring ciphertexts internally.\\\hline
\textbf{BFV / CKKS / BGV}&A relinearized ciphertext is $(A,B)\in R_q^2$. Multiplication temporarily produces three components.\\\hline
\end{tabularx}
\caption{\textbf{Ciphertext Representation}}
\end{table}

\begin{table}[htbp]
\small
\renewcommand{\arraystretch}{1.15}
\begin{tabularx}{\linewidth}{|p{0.17\linewidth}|X|}
\hline
\textbf{TFHE}&A finite message alphabet, often $m\in\mathbb Z_t$, embedded in the torus. Power-of-two $t$ is an encoding choice, not a universal requirement.\\\hline
\textbf{BFV / BGV}&$M\in R_t$. The modulus $t$ need not be prime. Full scalar batching commonly uses prime $t\equiv1\pmod{2n}$; other choices may yield extension-ring slots.\\\hline
\textbf{CKKS}&Up to $n/2$ complex slots, encoded approximately as an integer polynomial. Real/complex semantic values are not exact finite-ring elements.\\\hline
\end{tabularx}
\caption{\textbf{Plaintext and Slot Domain}}
\end{table}

\begin{table}[htbp]
\small
\renewcommand{\arraystretch}{1.15}
\begin{tabularx}{\linewidth}{|p{0.17\linewidth}|X|}
\hline
\textbf{TFHE}&Here $s_i\in\{0,1\}$. A separate GLWE secret is used in blind rotation.\\\hline
\textbf{BFV / CKKS / BGV}&A small polynomial $S\in R$, often with coefficients in $\{-1,0,1\}$, embedded into $R_q$. This is not multiplication in $R_3$.\\\hline
\end{tabularx}
\caption{\textbf{Secret Key}}
\end{table}

\begin{table}[htbp]
\small
\renewcommand{\arraystretch}{1.15}
\begin{tabularx}{\linewidth}{|p{0.17\linewidth}|X|}
\hline
\textbf{TFHE}&Fresh mask $\vec a\leftarrow\mathbb Z_q^k$; scalar $b\in\mathbb Z_q$. The simple integer encoding below assumes $t\mid q$.\\\hline
\textbf{BFV / CKKS / BGV}&Fresh mask $A\leftarrow R_q$; $B\in R_q$. An evaluated ciphertext need not have the distribution of a fresh encryption.\\\hline
\end{tabularx}
\caption{\textbf{Ciphertext and Fresh Randomness}}
\end{table}

\begin{table}[htbp]
\small
\renewcommand{\arraystretch}{1.15}
\begin{tabularx}{\linewidth}{|p{0.17\linewidth}|X|}
\hline
\textbf{All schemes}&Fresh error is sampled from a specified small distribution. Discrete Gaussian, bounded, and other distributions require their own parameter analysis.\\\hline
\textbf{Evaluation}&For independent centered errors, $\operatorname{Var}(e_1+e_2)=\operatorname{Var}(e_1)+\operatorname{Var}(e_2)$. Reusing one ciphertext gives error $2e$, not two independent samples.\\\hline
\end{tabularx}
\caption{\textbf{Noise}}
\end{table}

\begin{table}[htbp]
\small
\renewcommand{\arraystretch}{1.15}
\begin{tabularx}{\linewidth}{|p{0.17\linewidth}|X|}
\hline
\textbf{TFHE}&$\Delta=q/t$ in this integer encoding, with $t\mid q$.\\\hline
\textbf{BFV}&$\Delta=\lfloor q/t\rfloor$; write $r=q-t\Delta$. The remainder $r$ enters the correctness bounds.\\\hline
\textbf{CKKS}&Actual scale $\Delta$ is metadata and changes during multiplication/rescaling. Require the entire phase, including error, to fit the current modulus.\\\hline
\textbf{BGV}&Plaintext is unscaled; the error is a multiple of $t$. Writing $\Delta=t$ for the error factor does not identify it with the BFV/CKKS plaintext scale.\\\hline
\end{tabularx}
\caption{\textbf{Scaling Factor}}
\end{table}

\begin{table}[htbp]
\small
\renewcommand{\arraystretch}{1.15}
\begin{tabularx}{\linewidth}{|p{0.17\linewidth}|X|}
\hline
\textbf{TFHE}&$b=[\Delta m+e-\langle\vec a,\vec s\rangle]_q$. Sample a fresh mask and error for each encryption.\\\hline
\textbf{BFV}&$B=[\Delta M+E-AS]_q$.\\\hline
\textbf{CKKS}&$B=[\lfloor\Delta M\rceil+E_{\rm enc}-AS]_q$. Below $E$ includes encoding and encryption error.\\\hline
\textbf{BGV}&$B=[M+tE-AS]_q$.\\\hline
\end{tabularx}
\caption{\textbf{Symmetric Encryption}}
\end{table}

\begin{table}[htbp]
\small
\renewcommand{\arraystretch}{1.15}
\begin{tabularx}{\linewidth}{|p{0.17\linewidth}|X|}
\hline
\textbf{TFHE}&$\phi=b+\langle\vec a,\vec s\rangle\equiv\Delta m+e\pmod q$.\\\hline
\textbf{BFV / CKKS}&$\Phi=B+AS\equiv\Delta M+E\pmod q$.\\\hline
\textbf{BGV}&$\Phi=B+AS\equiv M+tE\pmod q$.\\\hline
\end{tabularx}
\caption{\textbf{Phase Relation}}
\end{table}

\begin{table}[htbp]
\small
\renewcommand{\arraystretch}{1.15}
\begin{tabularx}{\linewidth}{|p{0.17\linewidth}|X|}
\hline
\textbf{TFHE}&$m=\lfloor[\phi]_q/\Delta\rceil\bmod t$ if $|e|<\Delta/2$.\\\hline
\textbf{BFV}&$M=\lfloor(t/q)[\Phi]_q\rceil\bmod t$. For a chosen integer lift of $M$, a sufficient condition is $\|tE-rM\|_\infty<q/2$, where $r=q-t\Delta$.\\\hline
\textbf{CKKS}&Decode $[\Phi]_q/\Delta$ through the canonical embedding. If $\|\Delta M+E\|_\infty<q/2$, coefficient error is $E/\Delta$; slot error must be bounded after embedding.\\\hline
\textbf{BGV}&$M=[\Phi]_q\bmod t$. A sufficient no-wrap condition is $\|M+tE\|_\infty<q/2$.\\\hline
\end{tabularx}
\caption{\textbf{Decryption and Sufficient Correctness Conditions}}
\end{table}

\begin{table}[htbp]
\small
\renewcommand{\arraystretch}{1.15}
\begin{tabularx}{\linewidth}{|p{0.17\linewidth}|X|}
\hline
\textbf{TFHE / BFV}&A fixed modulus $q$ is sufficient for the basic presentation; RNS and modulus-switching variants also exist.\\\hline
\textbf{CKKS}&$q_i=\prod_{j=0}^{i}w_j$ with pairwise coprime factors. Dropping $w_l=q_l/q_{l-1}$ by rescaling changes scale by exactly $1/w_l$, even when $w_l\approx\Delta$.\\\hline
\textbf{BGV}&$q_i=\prod_{j=0}^{i}w_j$ with factors coprime to $t$. Taking every $w_j\equiv1\pmod t$ permits a simple residue-preserving convention; other chains require correction factors.\\\hline
\end{tabularx}
\caption{\textbf{Ciphertext Modulus}}
\end{table}

\begin{table}[htbp]
\small
\renewcommand{\arraystretch}{1.15}
\begin{tabularx}{\linewidth}{|p{0.17\linewidth}|X|}
\hline
\textbf{All schemes}&Add corresponding components: $(A_1+A_2,B_1+B_2)$, or $(\vec a_1+\vec a_2,b_1+b_2)$. Keys and moduli must match; CKKS scales must also be aligned.\\\hline
\textbf{Error}&The lifted phase errors add: $E_{\rm out}=E_1+E_2$. For BFV, re-expressing a wrapped sum as a representative in $R_t$ also introduces the $q-t\Delta$ carry correction.\\\hline
\end{tabularx}
\caption{\textbf{Ciphertext-to-Ciphertext Addition}}
\end{table}

\begin{table}[htbp]
\small
\renewcommand{\arraystretch}{1.15}
\begin{tabularx}{\linewidth}{|p{0.17\linewidth}|X|}
\hline
\textbf{TFHE / BFV}&Add $\Delta c$ to $b$, or $\Delta C$ to $B$, using the agreed integer representatives.\\\hline
\textbf{CKKS}&Add the integer encoding $\lfloor\Delta C\rceil$ to $B$; encoding error is included in the precision budget.\\\hline
\textbf{BGV}&Add $C$ to $B$ modulo the ciphertext modulus.\\\hline
\end{tabularx}
\caption{\textbf{Ciphertext-to-Plaintext Addition}}
\end{table}

\begin{table}[htbp]
\small
\renewcommand{\arraystretch}{1.15}
\begin{tabularx}{\linewidth}{|p{0.17\linewidth}|X|}
\hline
\textbf{TFHE / BFV / BGV}&Multiply both components by the unscaled plaintext $c$ or $C$. The phase error becomes $ce$ or $CE$, plus any representative correction. In $R$, $\|CE\|_\infty\leq\|C\|_1\|E\|_\infty$.\\\hline
\textbf{CKKS}&Multiply by an integer encoding at scale $\Delta_C$. If rescaling removes $w_l$, the output scale is $\Delta_{\rm out}=\Delta\Delta_C/w_l$. Component rounding contributes $\rho_B+\rho_A S$ with $\|\rho_A\|_\infty,\|\rho_B\|_\infty\leq1/2$.\\\hline
\end{tabularx}
\caption{\textbf{Ciphertext-to-Plaintext Multiplication}}
\end{table}

\begin{table}[htbp]
\small
\renewcommand{\arraystretch}{1.15}
\begin{tabularx}{\linewidth}{|p{0.17\linewidth}|X|}
\hline
\textbf{TFHE}&External products require a compatible GSW/GGSW encryption of one operand. Converting an LWE ciphertext to such a representation uses circuit bootstrapping, not ordinary PBS alone. Noise includes decomposition and evaluation-key errors.\\\hline
\textbf{BFV}&In the standard FV multiplication, form the integer negacyclic products $(D_2,D_1,D_0)=(A_1A_2,A_1B_2+A_2B_1,B_1B_2)$, scale components by $t/q$ and round, then relinearize modulo $q$. Do not reduce products modulo $q$ before scaling. The chapter also analyzes a separately specified $1/\Delta$ teaching variant.\\\hline
\textbf{CKKS}&Tensor multiplication, relinearization, and (when scheduled) rescaling. With input scales $\Delta_1,\Delta_2$, multiplication gives scale $\Delta_1\Delta_2$; after rescaling it is $\Delta_1\Delta_2/w_l$. Relinearization and rounding add error.\\\hline
\textbf{BGV}&Tensor multiplication and relinearization preserve the message modulo $t$ when the noise bound holds. Optional modulus switching uses the ratio $q_{l-1}/q_l$ and residue correction. It is not division by $t$ followed by rounding to a multiple of $t$.\\\hline
\end{tabularx}
\caption{\textbf{Ciphertext-to-Ciphertext Multiplication}}
\end{table}

\begin{table}[htbp]
\small
\renewcommand{\arraystretch}{1.15}
\begin{tabularx}{\linewidth}{|p{0.17\linewidth}|X|}
\hline
\textbf{TFHE}&Leveled operations have a finite noise budget. Gate/PBS evaluation performs bootstrapping and therefore does not belong in a claim about computation without bootstrapping.\\\hline
\textbf{BFV}&Fixed parameters permit a bounded circuit family determined by noise growth. A fixed ciphertext modulus does not imply unlimited multiplication.\\\hline
\textbf{CKKS / BGV}&The modulus chain and noise/precision budget constrain multiplicative depth and level consumption, not the total number of parallel multiplication gates. Refreshing requires a valid bootstrapping configuration.\\\hline
\end{tabularx}
\caption{\textbf{Computation Budget without Bootstrapping}}
\end{table}

\begin{table}[htbp]
\small
\renewcommand{\arraystretch}{1.15}
\begin{tabularx}{\linewidth}{|p{0.17\linewidth}|X|}
\hline
\textbf{Ring form}&Let $A=\sum_i d_i g_i$ be an exact gadget decomposition. Evaluation-key components $K_i$ have phase $g_iS+\varepsilon_i$ under $S'$. Then $c'=(0,B)+\sum_i d_iK_i$ has phase $B+AS+\sum_i d_i\varepsilon_i$ under $S'$.\\\hline
\textbf{Conditions}&TFHE uses the corresponding vector formula. Approximate decomposition adds a residual term. BGV evaluation-key errors must be multiples of $t$. A scalar added to a ciphertext means addition to its body component only.\\\hline
\end{tabularx}
\caption{\textbf{Key Switching}}
\end{table}

\begin{table}[htbp]
\small
\renewcommand{\arraystretch}{1.15}
\begin{tabularx}{\linewidth}{|p{0.17\linewidth}|X|}
\hline
\textbf{ModDrop}&For $q'\mid q$, reduce components modulo $q'$ without scaling. CKKS scale metadata stays unchanged; preserving the represented value requires the remaining phase to fit. BGV plaintext is preserved if $\|M+tE\|_\infty<q'/2$; arbitrary wraparound may change it modulo $t$.\\\hline
\textbf{BGV switch}&In the convention $q\equiv q'\pmod t$, choose each $c'\equiv c\pmod t$ close to $(q'/q)c$. Scaling changes absolute noise and modulus together, and adds rounding error; it does not sample a fresh noise distribution.\\\hline
\end{tabularx}
\caption{\textbf{Modulus Drop and Modulus Switching}}
\end{table}

\begin{table}[htbp]
\small
\renewcommand{\arraystretch}{1.15}
\begin{tabularx}{\linewidth}{|p{0.17\linewidth}|X|}
\hline
\textbf{TFHE}&Even a scalar requires an encoding such as $m\mapsto\Delta m$ and a decoding/noise margin. LUT padding and the valid input interval must be specified.\\\hline
\textbf{BFV / BGV}&Batch encoding is a polynomial CRT map over the plaintext ring, under the required factorization assumptions. It is optional when working directly with polynomial messages.\\\hline
\textbf{CKKS}&Canonical embedding and its inverse relate conjugate-symmetric complex evaluations to real coefficients; scaling and integer rounding then produce a finite encoding.\\\hline
\end{tabularx}
\caption{\textbf{Encoding and Decoding the Plaintext}}
\end{table}

\begin{table}[htbp]
\small
\renewcommand{\arraystretch}{1.15}
\begin{tabularx}{\linewidth}{|p{0.17\linewidth}|X|}
\hline
\textbf{Ring schemes}&Apply the automorphism $X\mapsto X^j$, then key-switch from $S(X^j)$ to $S(X)$. With the usual ordering, $j=5^h\bmod2n$ rotates CKKS complex slots or each of the two BFV/BGV rows; it is not one cycle through all $n$ scalar slots.\\\hline
\textbf{TFHE}&The scalar LWE interface has no SIMD slot rotation. Packed GLWE constructions have other polynomial operations and should be distinguished from this interface.\\\hline
\end{tabularx}
\caption{\textbf{Input Slot Rotation}}
\end{table}

\begin{table}[htbp]
\small
\renewcommand{\arraystretch}{1.15}
\begin{tabularx}{\linewidth}{|p{0.17\linewidth}|X|}
\hline
\textbf{TFHE / BFV / BGV}&Obtain an encryption of the same message (or a specified LUT output) with a usable new noise budget. The output still contains noise and may use another key or parameter set.\\\hline
\textbf{CKKS}&Refresh usable levels while controlling additional approximation error. Previously accumulated error in the represented value is not generally removed.\\\hline
\end{tabularx}
\caption{\textbf{Bootstrapping Goal}}
\end{table}

\begin{table}[htbp]
\small
\renewcommand{\arraystretch}{1.15}
\begin{tabularx}{\linewidth}{|p{0.17\linewidth}|X|}
\hline
\textbf{TFHE}&Modulus-switch to $2n$. With the plus-phase convention used here, the effective rotation error is $\hat e=(2n/q)e+\delta_b+\sum_i s_i\delta_{a_i}$. Initialize a GLWE accumulator encrypting $\Delta_{\rm out}V$, where $V$ is the unscaled LUT polynomial. Rotate first by $X^{-\hat b}$, then use encrypted secret bits to apply $X^{-\hat a_i s_i}$. Extract the constant term, which encodes $f(m)$ for the chosen LUT. The extraction key is derived from the GLWE secret; key-switch if needed.\\\hline
\textbf{TFHE conditions}&Define the unscaled lookup by $F(j)=\operatorname{coeff}_0(VX^{-j})$. It must be constant over the total error interval and respect negacyclic anti-periodicity $F(j+n)=-F(j)$. Arbitrary functions require a suitable restricted/padded input domain or an extended construction.\\\hline
\textbf{BFV}&For plaintext modulus $p^r$, transform decryption into rounding modulo $p^e$: $M=[\lfloor[c_0'+c_1' S]_{p^e}/p^{e-r}\rceil]_{p^r}$, under the modulus-switch error bound. Homomorphically evaluate the inner product using an encrypted secret; the constant body must be encoded at the new scale too. CoeffToSlot, digit extraction with exact divisions, and SlotToCoeff implement the remaining map. For odd $p$ use balanced digits; for $p=2$ include the half-unit rounding offset before deleting low bits. The CRT slot route requires $\gcd(p,2n)=1$; a binary bootstrap therefore needs a different compatible cyclotomic ring or construction. Their parameters and extension-slot layout are essential; see \cite{cryptoeprint:2022/1363}.\\\hline
\end{tabularx}
\caption{\textbf{Bootstrapping Details: TFHE, BFV}}
\end{table}

\begin{table}[htbp]
\small
\renewcommand{\arraystretch}{1.15}
\begin{tabularx}{\linewidth}{|p{0.17\linewidth}|X|}
\hline
\textbf{CKKS}&Modulus raising exposes a lifted phase $\Delta M+E+Kq_0$. CoeffToSlot is a homomorphic linear transform. EvalMod is a nonlinear approximation to modular reduction on a specified bounded input set. SlotToCoeff returns the result to coefficient form; all transforms consume budget and add error.\\\hline
\textbf{CKKS error}&The illustrative map $g(x)=q_0\sin(2\pi x/q_0)/(2\pi)$ satisfies $g(r+Kq_0)=r+\varepsilon_{\sin}$, with $|\varepsilon_{\sin}|\leq(2\pi)^2|r|^3/(6q_0^2)$. It is not an exact remainder operation; polynomial approximation and homomorphic evaluation add further errors.\\\hline
\textbf{BGV}&Transform the decryption relation, evaluate the secret-key inner product homomorphically, use CoeffToSlot and digit extraction to remove the unwanted digits, then apply SlotToCoeff. The CRT coprimality condition, odd-prime balanced digits, and binary rounding caveat are as in the BFV row. Exact division by $p$ changes plaintext modulus $p^j\to p^{j-1}$ and requires a divisibility invariant. Merely viewing a small-modulus ciphertext at a larger modulus does not perform this computation; see \cite{cryptoeprint:2022/1363}.\\\hline
\end{tabularx}
\caption{\textbf{Bootstrapping Details: CKKS, BGV}}
\end{table}

\begin{table}[htbp]
\small
\renewcommand{\arraystretch}{1.15}
\begin{tabularx}{\linewidth}{|p{0.17\linewidth}|X|}
\hline
\textbf{TFHE / BFV / BGV}&Track a bound on the complete phase error and the probability of exceeding the decoding margin. Bootstrapping creates a new bounded error; it does not set error to zero.\\\hline
\textbf{CKKS}&Track scale, magnitude, and slot error at each level. Increasing scale can reduce relative rounding error but consumes modulus capacity; it is not a cost-free or universal cure. Error need not increase monotonically in every circuit.\\\hline
\textbf{BGV}&Modulus switching reduces absolute phase magnitude while lowering the modulus and adding correction error. It is a noise-management step, not an independent refresh operation.\\\hline
\end{tabularx}
\caption{\textbf{Noise and Precision Management}}
\end{table}
